This section first introduces symbolic-execution-based specification generation for code segments without loops, and then presents the generation of general function specifications. Loop-invariant generation is deferred to the next section.

\subsection{Symbolic execution and strongest postcondition} %
Before stating the proposition, we fix its assumptions. 


At program point $p$, our \emph{flattened memory model} starts from in-scope variables and parameters and enumerates a finite set $\mathcal{L}_p$ of atomic cells reached by pointer dereferences, struct-field selection, and element-wise traversal of statically known or parameter-bounded arrays. Each scalar or pointer cell receives a typed symbolic variable. Access paths resolved to the same cell share a location; pairs resolved as distinct are marked separated. Values, aliasing, and updates are therefore expressed over $\mathcal{L}_p$.

The model produces only finite, first-order facts over named cells; it does not use recursive or inductive predicates. These facts serve as preconditions and postconditions for two back-ends: QCP (in QCPL) for symbolic execution, and Frama-C/WP (in ACSL) for verification. What limits the expressiveness of the facts is Frama-C/WP's requirement for automatic SMT discharge: because the verification step must be fully automatic, the facts are expressed in first-order logic without recursive or inductive predicates. \yfp{Hence the model is deliberately finite and handles only memory accesses whose target locations are statically enumerable. It supports scalars, structs, fixed arrays, and pointers with statically resolvable aliasing. Unbounded or recursive structures (e.g., linked lists and trees), dynamic unbounded allocation, and accesses with a symbolic offset (e.g., pointer arithmetic) fall outside the exact fragment and trigger the LLM fallback described in Section~\ref{sec:llm-fallback}.}





%Both back-ends consume the resulting precondition and postcondition facts---QCP in QCPL during symbolic execution and Frama-C/WP in ACSL during validation---but it is Frama-C/WP's automatic SMT discharge, not QCP's separation logic, that bounds the usable fragment: keeping the facts first-order and shape-predicate-free is precisely what lets Frama-C/WP discharge them directly, whereas QCP could in principle express more. The model is deliberately finite. It captures scalar, struct, fixed-extent array, and pointer state with statically resolvable aliasing, but it cannot name an unbounded or recursively defined set of cells; dynamically allocated or recursive heap shapes (e.g., linked lists and trees), unbounded allocation, and accesses whose target cannot be resolved to a finite location set (e.g., pointer arithmetic with a symbolic offset) therefore lie outside the exact fragment and trigger the LLM fallback of Section~\ref{sec:llm-fallback}.

% Within this model, the symbolic execution of statements maintains a symbolic state $(PC_\pi, \sigma_\pi)$ for each execution path $\pi$, where
% \begin{itemize}
%     \item $PC_\pi$ denotes the path condition, i.e., a logical formula over initial symbolic values of $\mathcal{L}_{p_0}$ and constants, that expresses the conditions required to traverse path $\pi$. $PC_\pi$ is updated when a branch is taken along $\pi$.
% \end{itemize}


Within this model, the symbolic execution of statements maintains a
\emph{symbolic state} $(PC_\pi, \sigma_{p,\pi})$ for each execution
path $\pi$.
At a program point $p$ on path $\pi$, let $\mathcal{L}^{\pi}_{p}$
denote the set of symbolic locations currently tracked by the
executor. This set includes the \yfp{function-entry locations}
$\mathcal{L}_{p_0}$ reachable from the function parameters,
together with the local variables \yfp{that are in scope at
$p$ along $\pi$} (hence
$\mathcal{L}_{p_0}\subseteq\mathcal{L}^{\pi}_{p}$).

\begin{itemize}
\item $PC_\pi$ denotes the path condition, i.e., a logical formula
over function-entry symbolic values and constants that encodes
the conditions needed to follow the path $\pi$.

\item $\sigma_{p,\pi}$ denotes the symbolic store at program
point $p$ along path $\pi$, i.e., a mapping from each currently
tracked symbolic location in $\mathcal{L}^{\pi}_{p}$ to a
symbolic expression over the function-entry values.
\yfp{Thus,
$\operatorname{dom}(\sigma_{p,\pi})=\mathcal{L}^{\pi}_{p}$.
We express the symbolic store $\sigma_{p,\pi}$ as the following
logical relation:
\[
\rho(\sigma_{p,\pi})
\;\equiv\;
\bigwedge_{\ell\in\mathcal{L}^{\pi}_{p}}
\ell^p=\sigma_{p,\pi}(\ell),
\]
where $\ell^p$ denotes the value of location $\ell$ at program
point $p$, and $\sigma_{p,\pi}(\ell)$ is a symbolic expression
over the function-entry values. Thus, $\rho(\sigma_{p,\pi})$
relates the function-entry state to the state at $p$ over
the tracked locations.}
\end{itemize}

The symbolic execution of a program is defined as a transformation over
sets of symbolic states, branching at conditional control flow points.
Formally, execution begins at the function entry point $p_0$ with the
initial symbolic state $(PC_0,\sigma_0)$, where $\sigma_0$ assigns an
initial symbolic value to each location in the entry point location set
$\mathcal{L}_{p_0}$, and $PC_0$ is the precondition constraining those
initial values. For a program point $p$, let the set of all symbolic
states reachable from $p_0$ be
$\{(PC_{\pi_i},\sigma_{p,\pi_i})\}_{i=1}^{m}$, where $\pi_i$ ranges over
all paths from $p_0$ to $p$.
\yfp{The pairs of function-entry states and reached states at $p$, represented over the tracked locations, satisfy
\[
\bigvee_{i=1}^{m}
\left(
PC_{\pi_i}\land\rho(\sigma_{p,\pi_i})
\right).
\]
%Here $\sigma_{p,\pi_i}(\ell)$ is an expression over the function-entry symbolic
%values, whereas $\ell^{p}$ denotes the corresponding value at program
%point $p$. An unchanged tracked cell retains the symbolic expression it had at the start of the current segment, and
%access paths known to alias refer to the same tracked location.
The exact claim applies only to the finite set of tracked locations $\mathcal{L}^{\pi}_{p}$, and 
memory locations outside this set are not constrained by the resulting formula. 
The path condition $PC_{\pi_i}$ is expressed over the same function-entry symbolic values and includes the starting
precondition. At the entry
point $p_0$, the store relation is the identity relation, so the
corresponding formula is $PC_0$ together with those identity
equalities.}


Given a segment $\mathcal F$, precondition $P$, and start point $p_{\mathit{start}}$, $\mathrm{SEStart}(\mathcal F,P,p_{\mathit{start}})$ explores feasible paths to the next boundary $p$ (a loop entry, call, or exit of $\mathcal F$) and returns $p$. The traversed segment contains no loop or call, so its input/output behavior can be summarized exactly.

We first define the notion of a \emph{strongest postcondition} in the Hoare‑logic sense for loop‑free segments.

\yfp{
\begin{definition}{(Strongest Postcondition)}
For a segment $\mathcal{F}$ and a precondition $P$, we consider postconditions whose truth, on each feasible path, depends only on the function-entry symbolic values and the endpoint values of locations tracked on that path. Within this scope, a formula $\Phi$ is the strongest postcondition of $\mathcal{F}$ with respect to $P$, if
\begin{itemize}
    \item $\{P\}\,\mathcal{F}\,\{\Phi\}$ holds: every execution under $P$ satisfies $\Phi$ at the segment endpoint;
    \item for any postcondition $\Psi$ in this scope such that $\{P\}\,\mathcal{F}\,\{\Psi\}$ holds, we have $\Phi\Rightarrow\Psi$.
\end{itemize}
\end{definition}
}


Let $\mathcal{F}$ be a segment from $p_{\textit{start}}$ to $p_{\textit{end}}$ with precondition $P$. \yfp{Symbolic execution resumes from the incoming symbolic state under $P$ at $p_{\textit{start}}$. This state relates the tracked values at the segment start to the function-entry symbolic values, which remain the basis of subsequent expressions. Let $\Pi=\textit{Path}(p_{\textit{start}},p_{\textit{end}})$ be the feasible paths explored under $P$, and $\Phi_{\mathit{SP}}=\textit{SP}(\mathcal{F},P,p_{\textit{start}},p_{\textit{end}})$. We define
\[
\Phi_{\mathit{SP}}
 = \bigvee_{\pi\in\Pi}\bigl(PC_\pi
 \land\rho(\sigma_{p_{\textit{end}},\pi})\bigr).
\]
Each disjunct uses the full symbolic store for its path; the tracked location sets may differ across paths.}

\yfp{
\begin{proposition}
\label{thm:sp-sound-complete}
$\textit{SP}(\mathcal{F},P,p_{\textit{start}},p_{\textit{end}})$ is the strongest postcondition of $\mathcal{F}$ with respect to $P$, if the following conditions hold:
\begin{itemize}
    \item[(A1)] Scope conditions: $\mathcal{F}$ contains no external calls (i.e., every called function has a definition), no loops, no recursive datatypes, and no recursive function calls. Every memory access is representable in the flattened model. For every input satisfying $P$, execution of the interval is well-defined under the modeled C semantics: memory accesses, aliasing, and arithmetic operations are represented exactly, and no undefined behavior occurs.
    \item[(A2)] Symbolic execution conditions: The incoming symbolic state exactly characterizes the segment-start states under $P$ and their relation to the function-entry values. Symbolic execution is path-complete: it enumerates all feasible paths from $p_{\textit{start}}$ to $p_{\textit{end}}$ under $P$. For each path $\pi$, $PC_\pi$ exactly characterizes the inputs that cause execution to follow $\pi$, and $\sigma_{p_{\textit{end}},\pi}$ exactly expresses the resulting values of all locations tracked on that path in terms of the function-entry symbolic values.
\end{itemize}
\end{proposition}
}

\yfp{The proof is given in \appref{app:sp-proof}.}

Unsupported features cause $\textit{SP}$ to return $\textit{SymExecFailed}$ and invoke the LLM fallback; \rev{undefined behavior instead produces an error.} Fallback specifications may be verifier-checked but do not inherit the strongest-postcondition guarantee of \yfp{Proposition~\ref{thm:sp-sound-complete}}.

\subsection{Default Precondition Generation}
\label{sec:default-pre}
\yfp{$\mathrm{DefaultPre}(\mathcal{F})$ generates sufficient preconditions for the supported symbolic-execution model. It emits three kinds of facts: validity (each pointer points to allocated memory), separation (distinct collected pointers point to initially disjoint memory regions), and bounds (restricting array lengths and integer values to the configured analysis ranges, e.g.\ $0 < \mathit{arr\_len} < \mathit{MAX\_LEN}$ and $0 \le v \le \mathit{MAX\_VAL}$ for integer value $v$). These restrictions define the input domain of the generated contract.

The generator processes the function parameters: for each pointer (including those inside structs and arrays), it emits a validity fact; after all pointers are collected, it appends a single separation fact asserting that their initial memory regions are pairwise disjoint. Validity and bounds help exclude common null-pointer and out-of-bounds runtime errors. \rev{QCP requires explicit initial memory predicates to justify dereferences, so the default separation fixes one layout rather than enumerating alias cases (\appref{app:qcp-memory}). Supported aliasing configurations can be specified by replacing the relevant default separation clauses with appropriate aliasing predicates. A function's actual specification therefore depends on its calling context and actual arguments.} The default generator does not emit additional functional assumptions such as sortedness.}



\rev{After any such replacement, additional designer-supplied precondition clauses are conjoined with the resulting precondition to form $PC_0$ at $p_0$ before symbolic execution begins.} Fig.~\ref{fig:example} illustrates both modes: the validity and separation facts on \texttt{CheckCalFun} are auto-generated, while the range constraints on its array values and length field are designer supplied.

In contrast, for recursive pointer-linked data structures (linked lists, trees, recursive structs), the number of reachable memory cells is determined only at runtime (by the list length or tree depth). Therefore, the static generator cannot emit a finite set of validity/separation facts for the entire shape; instead, the shape must be captured by an inductive predicate. In this case $\mathrm{DefaultPre}$ delegates the generation of the inductive shape predicate to the LLM fallback layer (Section~\ref{sec:llm-fallback}), which proposes the predicate and the corresponding precondition. The non-recursive part (validity of the head pointer and of primitive fields) is still emitted by the static generator. The LLM-generated predicate and precondition clauses are kept only if Frama-C/WP accepts them.

 
